% !TeX root = ../papers/nc-p3.tex
\section{平面に分解する分離三次式による有限切断の識別}
\label{sec:finite-split-separator-slices}
\begin{theorem}\label{thm:finite-split-separator-slices}
有限窓で全中央支持が4とする。ある局所点の分離三次式に二重平面因子があれば、
その反対切断は三次元の交代行列空間に一つの対称な階数一行列を加えた空間に左右同値である。
三次式は異なる二平面の積 $h s^2$、階数一切断は平方零の長さ3で、唯一の階数一点で $h\ne0$ である。
三重平面は生じない。分離三次式が三つの異なる平面に分かれるなら、その切断の階数一軌跡は
正次元または長さ2以下である。
\end{theorem}
\begin{proof}[二重因子]
写像 $E\to k^3$ の非零の左右評価はどれも階数2以上である。
二重因子の平面をなす三空間の行列 $M(x)$ は一般階数2である。
一般階数1以下なら共通左核または共通右核があり、この評価下界に反する。
同じ下界から一般階数2の $M$ にも定数左右核はない。
$\operatorname{adj}M$ は二次の階数1行列なので原始な列と行に因数分解し、
定数核がないため両ベクトルはともに一次である。
右核が三つの独立な一次式なら、核座標を $x$ として $M(x)x=0$ から
$M(x)=D[x]_\times$ である（各行の係数が交代）。$D$ は定数左核の不在から可逆。
交代空間に第四の行列 $J$ を加えて $\det([x]_\times+sJ)$ を取ると $s$ の一次係数は
$x^t Jx$ である。これが零なら $J$ も交代で四次元性に反する。左側も同様である。

従って両核はそれぞれ二つの独立な一次式を用いる。右核を $(u_2,-u_1,0)$ とすれば
最初の二列は $u_1,u_2$ と共通定数列の積である。その列を第一基底に取ると
下二行の最終列の二式 $v_1,v_2$ が独立である。二重因子条件は
$(v_2,-v_1)K(u_2,-u_1)^t=0$ であり、$K$ は第四行列の下二行・初め二列のブロックである。
$K$ は可逆である。特異なら下二行のある結合は全ての $E$ 上で最終列だけに値を持ち評価下界に反する。
もし二つの一次式空間が異なれば各基底変換で $(u_1,u_2)=(x,y),(v_1,v_2)=(x,z)$
にでき、四積 $x^2,xy,xz,yz$ の独立性から $K=0$ となる。よって同じ二空間である。
行列の基底を揃えると三空間は
$\bigl(\begin{smallmatrix}a&b^t\\ b&0\end{smallmatrix}\bigr)$ となる。残る一成分 $a$ が
$b$ の二座標と独立なのは三空間の単射性による。第四行列の外側成分を行・列加算で消せば
定理\ref{thm:finite-marked-planar-triple}の $E_D$ 型で $D$ は可逆である。
行列式 $a s^2\det D-s b^t\operatorname{adj}(D)b$ が $s^2$ で割れるので $D$ は交代。
第一行または第一列の符号変更で表示した交代空間と階数一行列の和になる。
階数一点とそのスキームも同定理の計算通りである。

\end{proof}
\begin{proof}[三つの異なる平面]
有限な長さ3以上の切断の全模型でこの分解が起きないことを確認する。
非Gorenstein三重点の場合、上の $a s\det D-b^t\operatorname{adj}(D)b$ は後項が非零なら既約二次式である
（$a$ について一次で、係数と定数項は互いに素）。後項が零なら二重因子の場合である。
Gorenstein三重点の三空間は長さ3環 $B$ の全乗法行列で、$E=m_B+kJ$ である。
行列式が一次因子に分かれるなら、単位行列を使った特性多項式の根は
$\chi_i(b)+s\lambda_i$ である。実際 $s=0$ で乗法行列の三指標に制限され、
各一次因子の $s$ 係数は定数である。
$B=k^3$ なら対角部分を除いた $J$ を選ぶ。各 $\lambda_i=0$ は行列式の $s$ 一次項より従う。
二次・三次項から $J_{ij}J_{ji}=0$、$\det J=0$ である。
有向グラフには二閉路がなく三閉路もない（三閉路の逆の積は零）。
同時に上三角化でき、共通不変直線が評価下界に反する。
$B=k[t]/t^3$ では $N=E_{12}+E_{23}$ として、$N,N^2$ の方向はどの固有値も変えない。
平方・三乗の跡の係数から $\operatorname{tr}(N^2J)=\operatorname{tr}(NJ)=
\operatorname{tr}(N^2J^2)=0$。従って $J_{31}=0,J_{21}=-J_{32},J_{21}J_{32}=0$ で共通不変直線がある。
$B=k[t]/t^2\times k$ では同様に $E_{12}$ を使い、$J_{21}=J_{23}J_{31}=0$ なので
共通不変直線または双対の共通不変直線がある。

長さ4の跡模型では、三行の行列式は非零定数倍を除き
$\tau(\det(m_a)a^{-1})$ である。実際跡超平面への形式の制限の行列式は、
周囲の行列式とその逆形式による法線 $\tau$ の自己対合との積であり、単元開集合上でこの式になる。
曲線状の長さ $d_i$ の因子では、剰余座標を $s_i$ としてノルムは $\prod s_i^{d_i}$。
$d_i\ge2$ の因子があればその最高冪の係数座標についてこの三次式は一次で、係数は
非零定数倍の $s_i^{d_i-2}\prod_{j\ne i}s_j^{d_j}$ である。
定数項はこの係数のどの一次因子でも割れない。実際 $s_j=0$ で第 $j$ 因子の
逆数分子の最高冪成分が残り、その最初の接座標の $d_j-1$ 乗の係数は完全跡のため非零である。
単純因子ではこの冪は0乗である。$j=i$ を調べる必要のある $d_i>2$ でも同じである。
ゆえに三次式は既約である。四つの単純因子の場合は
$\sum c_i\prod_{j\ne i}s_j$ ($c_i\ne0$) を $s_1$ の一次式として同じ互いに素の判定を使う。
局所 $(1,2,1)$ の場合は $s$ と、$sw$ 項及び二つの接変数の非退化二次形式を持つ
非退化二次式との積である（$w$ はsocle係数）。以上で三つの異なる一次因子を排除した。
\end{proof}

\begin{corollary}\label{cor:single-support-quartic-slice}
有限な共通切断が長さ4で、点スキームの台が一点なら、任意の接次元の印で
行・列側の反対対称性が成立する。反対階数8で台が二点以下の残余候補でも、
この切断からの非対称な外積の支持環が曲線状の長さ4である場合は排除される。
\end{corollary}
\begin{proof}
接次元2だけを扱えばよい。定理\ref{thm:finite-artin-kernel-symmetry}の外積の
支持環 $S$ と三次持上げを使う。$S$ の一つの局所socleに対応する三次分離式は
曲線状因子の積なら $H\prod_i l_i^{d_i}/l_j$ である。ここで $d_i$ は各局所長、
$l_i$ は各方向での $L$ の剰余一次式、$j$ はそのsocleの因子である。
これは $C=H c N L^{-1}$ を第 $j$ の指標で読むことで従う。
socle元を $\eta$ 倍しても元の点環の非零socle元になるので、
定理\ref{thm:finite-cubic-separators}の当該点の行列式とこの三次式は定数倍で一致する。
局所 $(1,2,1)$ の $S$ では同様に分離式は $H s^2$ となる。

台が一点で曲線状 $S$ の長さが4なら分離式は三重平面となり不可能である。
残る長さ3または局所 $(1,2,1)$ 型では $H$ は一次で、分離式は $H s^2$ である。
上の定理から $H$ は印で非零なので点環全体で単元である。
しかし $H$ と三次持上げの式から、逆数分子の非零な二次座標（例えば $s^2$）は
$Z=\operatorname{Spec}S$ に支持される。隣の偶座標は $S$ に全射するのでその平方も全射し、
平方どうしの完全跡直交によりこの二次座標は零になる。これは二次式の単射性に反する。
全射の場合および共線因子の排除は既出の通りで、奇側の対称性が従う。

台が二点で $S$ が曲線状長さ4かつ局所なら同じ三重平面の矛盾となる。
二因子なら $Z$ は印を含む。$H$ は定数であり、印でのsocle分離式は
印で消える三つの一次式の積なので射影三空間で印における重複度が3となる。
一方、階数一点での行列式曲面の重複度は正確に2である。
その点を $E_{11}$ としたとき二次初期式は右下二行二列の行列式で、恒等零なら
そのブロックの線型空間は共通左核か共通右核を持つ。従って全四空間の一側評価が
高々一次元となり、純関係の不在に反する。これで二因子の場合も除かれる。
\end{proof}
